% =============================================================================
% Discussion section excerpt -- covers EC50,L, Emax,L, k_immune, and the
% compartment exchange rates. This is NOT a full paper: it holds only the
% Discussion section so it can be compiled and reviewed standalone. Merge
% \section{Discussion} below into the main manuscript .tex file; the main
% file's own preamble must load amsmath and graphicx (both used here).
% =============================================================================
\documentclass[11pt]{article}

\usepackage[margin=1in]{geometry}
\usepackage{amsmath,amssymb}
\usepackage{graphicx}
\usepackage[colorlinks=true,linkcolor=blue,citecolor=blue]{hyperref}

\title{Discussion (excerpt)}
\author{}
\date{}

\begin{document}
\maketitle

\section{Discussion}

Across four mechanistically unrelated parameters -- linezolid potency
($EC_{50,L}$), linezolid's bacteriostatic ceiling ($E_{max,L}$), host
immune clearance ($k_{immune}$), and the rates of exchange between blood and
reservoir ($f_{r \to b}$, $f_{b \to r}$) -- the model returns the same
qualitative result: the blood ($R_b$) and reservoir ($R_{res}$) resistant
subpopulations never escape or clear independently. In every deterministic
sweep and every Monte Carlo sample, the two compartments cross the limit of
detection together, defining a single shared boundary between a clearing and
a relapsing infection rather than two separate ones. That consistency,
reproduced across four independent analyses, is itself a principal finding
of this work, and it shapes how each individual parameter should be
interpreted below.

Linezolid potency produces the least intuitive result of the four
parameters, because the model's calibrated baseline itself, not just its
sampled variability, is the striking finding. $EC_{50,L}$ is held fixed at
$1.0~\mathrm{mg/L}$ in every other analysis in this paper, and that same
value is the arithmetic mean of the log-normal distribution sampled here --
yet $1.0~\mathrm{mg/L}$ sits \emph{above} the bisected escape boundary
($EC_{50,L} \approx 0.93~\mathrm{mg/L}$). In other words, the model's own
calibrated operating point for linezolid potency already predicts a
relapsing, not a clearing, infection, before any inter-patient variability
is introduced at all -- the same pattern seen below for $k_{immune}$. The
reported median ($0.70~\mathrm{mg/L}$) is a different quantity, pulled well
below that mean by the heavy right skew of a $\mathrm{CV}=1.117$ log-normal,
and it happens to fall on the clearing side of the threshold. The practical
result is a reversal of the naive expectation: it is not that variability
introduces a one-in-three chance of relapse into an otherwise safe baseline;
it is that the deterministic baseline alone predicts relapse, and it is the
variability -- specifically, the heavy tail of a right-skewed distribution
pulling most patients' effective potency below the population mean -- that
gives 65.3\% of simulated patients a chance at clearance despite operating,
on average, at the same potency as the model's own escaping baseline.
Clinically, the underlying mechanism is not exotic: sub-MIC drug exposure
from poor tissue penetration, protein binding, or early low-grade
heteroresistance could all plausibly shift an isolate's effective
$EC_{50,L}$, and the true population distribution of such effects is
unlikely to be symmetric around its mean.

$E_{max,L}$ is the more consequential of the two pharmacodynamic parameters
examined. Because $E_{max,L}$ is a ceiling rather than a potency, escape
occurs \emph{below} the threshold: a weaker maximum bacteriostatic effect,
not a right-shifted dose-response curve, is what allows resistant growth to
outrun drug pressure. The bisected threshold
($0.8161~\mathrm{h}^{-1}$) sits just above the assumed sampling mean
($0.8~\mathrm{h}^{-1}$), so more than half of the simulated population
(56.5\%) draws an $E_{max,L}$ value insufficient to clear the infection.
Unlike $EC_{50,L}$, where escape is a minority outcome under realistic
variability, $E_{max,L}$ variability alone is sufficient to make relapse the
\emph{majority} outcome. This asymmetry is worth emphasizing to a clinical
audience: potency (captured by MIC-type susceptibility testing) is the
parameter clinicians already track, while the ceiling effect -- which
depends on local drug exposure, tissue penetration, and any biofilm- or
persister-associated tolerance in the bone reservoir -- is not routinely
measured at all, despite the model identifying it as the more fragile of
the two.

Host immune clearance is the one parameter where the model's own calibrated
default is not a robust operating point. The bisected escape threshold,
$k_{immune} \approx 0.15~\mathrm{h}^{-1}$, sits \emph{above} the model's own
baseline assumption ($k_{immune} = 0.12~\mathrm{h}^{-1}$), meaning the
default patient in this model already lacks sufficient endogenous immune
clearance to resolve the resistant subpopulations without antibiotic
pressure -- consistent with the underlying clinical picture of a
reservoir infection (osteomyelitis progressing to endocarditis) that host
immunity alone has already failed to control. Sampling $k_{immune}$ from a
distribution centered above this threshold (mean $0.17~\mathrm{h}^{-1}$,
$\mathrm{CV}=0.25$) still yields clearance in most simulated patients
(64.3\%), but the remaining 35.7\% relapse from immune variability alone.
Clinically, this maps onto exactly the comorbidities known to predict poor
outcomes in osteoarticular \textit{S. aureus} infection -- diabetes,
immunosuppression, advanced age -- and suggests that host factors are not a
secondary modifier of treatment response in this model but a comparably
large driver of relapse risk to the drug's own pharmacodynamics.

The two exchange rates behave differently from the three pharmacodynamic/
host parameters above, both in method (a deterministic log-spaced grid
rather than log-normal Monte Carlo sampling) and in structure: rather than
a single threshold separating escape from clearance, each rate exhibits two
flat, insensitive regimes bracketing one narrow transition band. That
structure has a direct mechanistic explanation once the exchange rate is
read as a timescale rather than a bare rate constant: $1/f$ is roughly how
long exchange takes to matter, and the transition band identified for both
rates, $[10^{-3}, 10^{-2}]~\mathrm{h}^{-1}$, corresponds almost exactly to
$1/(10^{-3}~\mathrm{h}^{-1}) \approx 42$~days and
$1/(10^{-2}~\mathrm{h}^{-1}) \approx 4$~days -- the linezolid and
vancomycin course lengths, respectively. Below $10^{-3}~\mathrm{h}^{-1}$,
exchange is too slow to complete even once during the entire linezolid
course, so the compartments behave as though decoupled and the outcome is
insensitive to the exact rate; above $10^{-2}~\mathrm{h}^{-1}$, exchange is
faster than even the shorter vancomycin window, so the two compartments
equilibrate and the outcome again saturates, this time at the well-mixed
limit. The rate only matters clinically in the narrow band where its own
timescale is comparable to a treatment-phase duration -- neither negligible
nor effectively instantaneous. This gives a general, testable rule for
compartmental models of this kind: a coupling parameter is worth measuring
precisely only when its reciprocal falls within the range of the therapy's
own phase durations, and can otherwise be treated as a coarse order-of-
magnitude estimate without materially affecting predicted outcomes. Given
how difficult exchange/penetration rates are to measure directly in bone
and cardiac vegetations \emph{in vivo}, this is a practically useful
distinction: it tells the modeler, and by extension the clinician
interpreting the model, which unmeasured parameters actually need to be
pinned down.

A further pattern emerges when the three Monte Carlo analyses are placed
side by side. Approximating each escape fraction as the probability that a
log-normal draw exceeds (or, for $E_{max,L}$ and $k_{immune}$, falls below)
the bisected threshold reproduces all three reported fractions to within a
few percentage points using only the threshold, the assumed median, and the
sampling $\mathrm{CV}$: $EC_{50,L}$ ($z \approx 0.32$, predicting
$\approx$37\% against a reported 34.7\%), $E_{max,L}$ ($z \approx 0.21$,
predicting $\approx$58\% against a reported 56.5\%), and $k_{immune}$
($z \approx -0.39$, predicting $\approx$35\% against a reported 35.7\%).
Every threshold, in other words, sits within roughly four-tenths of a
standard deviation of its own sampled median. This does not mean the Monte
Carlo results are circular -- the threshold itself is located independently
by deterministic
bisection on the mechanistic model, not fit to the sampling distribution --
but it does mean the escape fractions above are better read as a
consequence of \emph{where the calibrated model happens to sit relative to
its own threshold} than as an independent measurement of clinical risk in
absolute terms. The finding worth stating plainly is that the calibrated
model is not comfortably clear of relapse in any of the three dimensions
examined; it is poised close to criticality in all three, independently.
Whether that reflects a property of this particular calibration or a more
general feature of relapsing reservoir infections that have already reached
the point of requiring salvage linezolid therapy is a question this model
cannot resolve on its own, but the consistency across three unrelated
parameters is at least suggestive of the latter.

Finally, the fact that $R_b$ and $R_{res}$ never escape independently
across any of the four analyses is itself a structural property of the
model worth making explicit rather than treating as an incidental
observation. The two resistant subpopulations are coupled by strictly
positive bidirectional exchange ($f_{r \to b}, f_{b \to r} > 0$) in every
parameter regime examined, which is a plausible mechanistic reason for the
two compartments to share fate rather than fail independently: near the
detection floor, growth in either compartment is damped by the model's
smooth-floor terms, but the linear exchange coupling between them is not,
so whichever compartment begins to escape first can reseed the other before
either clears. A fully rigorous stability argument for this coupling --
accounting for the model's nonlinear floor terms rather than a naive
linearization at exact extinction -- is beyond the scope of the present
analysis, but formalizing it is a natural next step, and would convert an
empirical regularity observed four times over into a proven structural
property of the model.

\end{document}
